\begin{abstract}
Рост нагрузки может замедлять отдельные операции и увеличивать повторную обработку, из-за чего перегрузка сохраняется после всплеска поступлений. Мы исследуем, какие зависимости должна сохранять исполнимая модель процесса для подготовки реакции по журналам, записанным до затяжной тяжёлой перегрузки. Метод обучает модели длительности и маршрутизации, сравнивает допустимые действия при заданных стрессах и фиксирует реакцию с наблюдаемым правилом активации до независимой проверки. В исходном последовательном эксперименте сохранение нагрузочных зависимостей снижает стоимость незавершённой работы и вмешательства на 38.2\% относительно статического планирования и на 35.4\% относительно предиктивного отбора компонентов. Независимое расширение охватывает 288 обучающих наборов и шесть структур с параллельными операциями, альтернативными маршрутами, несколькими возвратами и конкурирующими ресурсными ограничениями. При нагрузочных возвратах сравнение с равными весами процессов снижает нормированную на отсутствие действий стоимость на 0.090 при заданной структуре; после автоматического обнаружения структуры общий выигрыш исчезает. Улучшения отдельных структур не проходят поправку на шесть основных сравнений, а один процесс с общим ресурсом ухудшается. Результат выявляет проблему моделирования для решений: зависимости с малой предиктивной ценностью до кризиса способны менять предпочтительную реакцию, тогда как структурные ошибки и необоснованная экстраполяция могут устранять пользу. Предиктивная точность, корректность структуры и фактическое качество подготовленной реакции требуют раздельной оценки.
\keywords{Выявление процессов \and Симуляция бизнес-процессов \and Перегрузка \and Подготовка реакции \and Зависимость от нагрузки}
\end{abstract}

\section{Введение}
В сложной информационной системе нагрузка может менять объём работы, необходимой для завершения каждого экземпляра. Отдельные операции замедляются, повторная обработка учащается. Оба эффекта дольше занимают ресурсы и усиливают накопление работы. Поэтому временный всплеск способен оставить систему перегруженной после возвращения поступлений к обычной интенсивности. Подготовка реакции требует понимания этих обратных связей наряду с ограничениями мощности.

Разные действия меняют разные части механизма. Дополнительный работник увеличивает мощность, сокращение возвратов уменьшает работу на один экземпляр. В процессно-ориентированной системе инженер должен выбрать допустимое действие, задать наблюдаемый порог и оценить системные последствия до внедрения. Параллельные ветви, повторяемые подпроцессы и общие ресурсы делают эти последствия зависимыми не только от локальной очереди, где проблема проявилась первой.

Предиктивная точность обычной эксплуатации не отвечает непосредственно на этот вопрос. Нагрузочная зависимость может почти не улучшать прогноз в наблюдавшемся диапазоне, но менять ранжирование действий при перегрузке. Её удаление способно сохранить приемлемое описание знакомых условий, одновременно ухудшив подготовленную реакцию. Обратно, сохранение слабо оценённой зависимости может породить ошибочную экстраполяцию. Поэтому выбор состава модели необходимо оценивать через его влияние на подготовленное действие.

Симуляция бизнес-процессов и предписывающий мониторинг дают основу для обучения моделей по журналам и оценки вмешательств~\cite{simod,alarms,hybridpolicy,rlpolicies}; нагрузочные длительности уже используются в гибридной симуляции~\cite{deepsimulator,hybridpolicy}. Мы рассматриваем три связанные задачи: сохранение зависимостей, существенных для выбора реакции, разделение предиктивного качества и качества решения, а также устойчивость выбора к сложной структуре и ошибкам её обнаружения.

Процедура восстанавливает наблюдаемое состояние процесса, обучает модели длительности и маршрутизации и сравнивает действия в заданной инженером библиотеке стрессов. Ресурсные ограничения, допустимые эффекты и стоимости задаются извне. Действие и порог WIP фиксируются до независимого тестирования. Сравнения сохраняют одинаковый поиск по политикам и набор действий, варьируя отбор поведенческих компонентов и структуру процесса; простая эвристика и планировщик с известным генератором помогают разделить вклад обученной модели и поиска.

Последовательный контрольный случай показывает, что предиктивный отбор способен удалить существенную зависимость возвратов. Структурное расширение проверяет этот результат на шести сконструированных классах процессов, отдельно с известной и обнаруженной структурой. Неоднородные результаты задают границы вывода: сохранение нагрузочных зависимостей может улучшать подготовку, но не заменяет проверку системной модели и не доказывает постоянного преимущества симуляции перед простым правилом. Инженерный вклад — исполнимая процедура подготовки и доказательства того, как выбор модели отражается на полученной реакции.

\section{Связанные работы}
Анализ очередей по журналам событий позволяет оценивать время ожидания~\cite{queue}. Исследования обслуживания при переменной нагрузке показывают, что может меняться и сама длительность обработки~\cite{load}. Работы о метастабильных отказах описывают, как временное возмущение запускает самоподдерживающийся рост объёма работы~\cite{metastable}. Мы изучаем эти механизмы через перегрузку процесса на конечном горизонте и доступные действия для её ограничения.

Выявление процессов и симуляция по данным дают исполняемые модели поведения системы. Inductive Miner и PM4Py восстанавливают структуру процесса~\cite{inductive,pm4py}; SIMOD выявляет и калибрует имитационные модели~\cite{simod}; Prosimos учитывает различия в поведении ресурсов~\cite{prosimos}; AgentSimulator строит модель вокруг исполнителей~\cite{agents}. DeepSimulator объединяет выявленную структуру, обученные модели времени и признаки взаимодействия задач~\cite{deepsimulator}. Эти подходы составляют основу моделирования в нашей работе. Наше сравнение касается информации о поведении, необходимой модели для выбора реакции при перегрузке.

Различие между ошибкой прогноза и качеством решения уже изучается в predict-then-optimize и decision-focused learning. В Smart Predict-then-Optimize ошибки прогноза оцениваются через потери последующей оптимизации~\cite{spo}; decision-focused learning обучает предиктивную модель с учётом целевой функции решения~\cite{dflsurvey}. Наш вклад не состоит в новой функции такого обучения. Мы исследуем сохранение компонентов исполнимой модели процесса, обученной до затяжного кризиса. Здесь нагрузка формируется внутри системы: вмешательство меняет очереди, очереди влияют на обслуживание и повторную обработку, а эти изменения определяют дальнейшую нагрузку. Исторические журналы содержат наблюдения обычного исполнения, а не оценки стоимости альтернативных реакций на кризис. Инженерная задача состоит в сохранении зависимостей, необходимых для сравнения таких реакций с учётом маршрутизации, синхронизации и общих ресурсов.

Ближайшие работы по обучению политик включают распределение ресурсов с учётом нагрузки. Meneghello и соавторы~\cite{hybridpolicy} объединяют PPO с гибридной симуляцией RIMS. Модель длительности использует число незавершённых задач и занятость ресурсов, маршрутизация задаётся вероятностями переходов, а оценка включает перегрузку из-за ограничений ресурсов. Middelhuis и соавторы~\cite{resourcepolicy} обучают назначение ресурсов с помощью глубокого обучения с подкреплением и приближения ценности на основе оценок, используя очереди и текущие назначения как состояние. Branchi и соавторы~\cite{rlpolicies} обучают выбор операций по журналам событий с помощью явных марковских процессов принятия решений и офлайн-обучения с подкреплением. Сравнение фиксированных вероятностей переходов с решениями по признакам и истории задачи показывает значение представления среды оценки.

В этом уже изучаемом контексте наша работа выделяет конкретный выбор представления: сохраняет ли отбор по точности прогноза до затяжного кризиса зависимости длительности и возвратов, которые становятся существенными под нагрузкой. Накопленная работа меняет обработку и повторные операции; эти изменения влияют на дальнейшее накопление работы. Мы сравниваем модели, сохраняющие или исключающие такие зависимости, при одинаковом каталоге действий и поиске по симуляциям.

Методы выбора вмешательств решают смежные задачи подготовки реакции. Интерпретируемые политики предлагают настраиваемые фильтрацию и ранжирование с оптимизацией параметров при ограниченных ресурсах~\cite{whitebox}; системы сигналов оптимизируют время вмешательства~\cite{alarms}; причинные методы оценивают эффект действий по историческим данным об их применении~\cite{causal}. Многокритериальные методы учитывают компромиссы между организационными показателями~\cite{multikpi}. Мы принимаем эффекты действий как инженерные входы и оцениваем стоимость накопленной работы и частоту затяжной перегрузки. Подходы к оценке качества симуляции различают точность структуры, времени и перегрузки~\cite{validation}; наша оценка добавляет фактическое последствие действия, выбранного моделью.
